% !TeX spellcheck = en_US
\documentclass[french]{yLectureNote}

\title{Électromagnétisme 3}
\subtitle{Physique}
\author{Paulhenry Saux}
\date{\today}
\yLanguage{Français}
%jesse.groenen@cemes.fr
\professor{M.Brut}
\usepackage{graphicx}%-pour mettre des images
\usepackage[utf8]{inputenc}%encodage
\usepackage{geometry}%pour modifier les tailles et mettre a4paper
%\usepackage{awesomebox}%pour les boites d'exercices, de pbq et de croquis d\'esactiv\'e pour les TP de PC
\usepackage{tikz}%pour deiffner + d\'ependance de chemfig
% \usepackage{tabularx}%pour dimensionner automatiquement les tableaux avec variable X
\usepackage{awesomebox}%Pour les boites info, danger et autres
\usepackage{menukeys}%Pour deiffner les touches de Calculatrice
\usepackage{fancyhdr}%pour les en-t\^ete personnalis\'ees
\usepackage{blindtext}%pour les liens
\usepackage{hyperref}%pour les liens (\`a mettre en dernier)
\usepackage{caption}%pour la francisation de la l\'egende table vers Tableau
\usepackage{pifont}
\usepackage{array}%pour les tableaux
\usepackage{lipsum}
\usepackage{yFlatTable}
\usepackage{multicol}
\usepackage{tabularx}
\usepackage{booktabs}
\usepackage{esint}
\newcommand{\Lim}[1]{\lim\limits_{\substack{1}}\:}
\renewcommand{\vec}{\overrightarrow}
\newcommand{\N}[0]{\mathbb{N}}
\newcommand{\dd}{\mathrm{d}}
\newcommand{\norm}[1]{||\vec{1}||}
\newcommand{\fo}{\psi(\vec{r},t)}
\newcommand{\HH}{\hat{H}}
\newcommand{\hb}{\hbar}
\newcommand{\lap}{\nabla^2}
\newcommand{\lapcc}{\frac{\partial^2 }{\partial x^2}+\frac{\partial^2 }{\partial y^2}+\frac{\partial^2 }{\partial z^2}}
\newcommand{\E}{\vec{E}(M)}
\newcommand{\B}{\vec{B}(M)}
\newcommand{\V}{V(M)}
\newcommand{\er}{\vec{e_{\rho}}}
\newcommand{\ez}{\vec{e_{z}}}
\newcommand{\ep}{\vec{e_{\varphi}}}
\newcommand{\pip}{\pi^+}
\newcommand{\pim}{\pi^-}
\newcommand{\mv}{\mathcal{V}}
\DeclareMathOperator{\Div}{Div}
\newcommand{\rot}{\vec{\mathrm{rot}}}
\newcommand{\grad}{\vec{\mathrm{grad}}}
\newcommand{\vds}{\vec{\dd S}}
\newcommand{\Ea}{\vec{E_{aux}}}
%\setcounter{chapter}{3}
\begin{document}

\chapter{Propagation d'onde dans la matière}
\section{Notions de cours}
Il n'y a pas d'absorption, donc $\underline{\epsilon_r} = \epsilon_r$ réel.

la fonction diléctrique relative correspond  à la permitivité relative.

Dans notre cas, $\omega << \omega_{domaines absorption}$ donc on est statique.

$\epsilon_s$ est statique, de plus $\epsilon_s >1$.
\subsection{Équations de Maxwell dans un milieu localement neutre}
On écrit E et B sous la forme 
$$\underline{\vec{E}} = \underline{\vec{E_m}}e^{-i\omega t}$$
On a donc dans un milieu localement neutre :
\begin{itemize}
    \item $\rot \underline{\vec{H}} = - \frac{\partial \underline{\vec{D}}}{\partial t}=-i\omega \underline{\vec{D}}$
    \item $\rot \underline{\vec{E}}=-\frac{\partial \underline{\vec{B}}}{\partial t} = i\omega \underline{\vec{B}}$
    \item $\Div \underline{\vec{D}} = 0$
    \item $\Div \underline{\vec{B}} = 0$
\end{itemize}
On a de plus 2 remations constitutives du milieu :
\begin{itemize}
    \item $\underline{\vec{D}} = \epsilon\underline{\vec{E}}$
    \item $\underline{\vec{H}} = \frac{\underline{\vec{B}}}{\mu_0}$
\end{itemize}
\subsection{Équation de propagation}
On utilise les équations 1 et 3 pour obtenir l'équation de propagation.

On utilise pour cela l'identité : $\rot\rot \underline{\vec{E}} = \grad(\Div{\underline{\vec{E}}})-\lap\vec{E}$

Donc :
\begin{flalign*}
   \rot \rot \underline{\vec{E}}&= i\omega \rot\underline{\vec{B}}\\
   & = i\omega \mu_0 \rot \underline{\vec{H}}\\
   &= \omega^2 \mu_0 \underline{\vec{D}}\\
   &= \omega^2\epsilon_0 \mu_0\epsilon_s \underline{\vec{E}}\\
   -\lap\underline{\vec{E}} &= 
   \Rightarrow \lap\underline{\vec{E}} + \frac{\omega^2}{c^2}\epsilon_s \underline{\vec{E}} = \vec{0}
\end{flalign*}
\subsection{Solution}
On cherche la solution sous la forme d'une onde plane : $\underline{\vec{E}} = \underline{\vec{E_m}} e^{-i\omega t}$ avec $\underline{\vec{E_m}} = \vec{E_0}e^{i\vec{k}\cdot \vec{r}}$

Ici, $\vec{K} = \underline{\vec{k}\cdot \vec{e_z}}$ donc $\vec{E_0} = E_0\vec{e_x}$%TODO POURQUOI c'est plus complexe

En injectant cette solutions dans l'équation de propagation :
$$\lap \underline{\vec{E}} = -k^2 e^{-i\omega t} \underline{\vec{E_m}}$$
 
% On obtient :
%  $$k^2 \vec{\underline{E}} = \frac{\omega^2}{c^2}\epsilon_r(\omega)\vec{\underline{E}}$$
 On peut en déduire la relation de dispersion du milieu :
 $$\omega^2 = \frac{c^2}{\epsilon_0^2}\underline{k^2}$$
\subsubsection{Parties imaginaires et complexes de n}
On a la relation 
$$\underline{\vec{n}}^2 = (n+i\kappa)^2 \underline{\epsilon_r} = \underline{\epsilon_s}$$
Donc 
$$n^2-\kappa² +2in\kappa = \epsilon_s$$
On en déduit que
$$n^2-\kappa^2 = \epsilon_s, 2n\kappa=0$$
Comme $\epsilon_s >1>0$ donc on choisit $n=\epsilon_s^{1/2}, \kappa=0$

On déduit de la relation de propagation que $\underline{k} = \pm \frac{n\omega}{c}>0$

Donc l'onde est de la forme $\vec{E_0}e^{i(\frac{\omega}{c}nz-\omega t)}$

L'amlitude réelle est constante, il y a  donc propagation de l'onde dans atténuation.

\subsection{Vitesse de phase}
La phase vaut $kz-\omega t$ et par définition, $v_{\varphi} = \frac{\omega}{k} = \frac{c}{\sqrt{\epsilon_s}}<c$
\subsection{Période spatiale}
On a $\lambda = \frac{2\pi}{k} = \frac{2\pi}{\omega}\frac{c}{n}$

Dans le vide, $\omega  = k_0c\Rightarrow \lambda_0 =\frac{2\pi c}{\omega}$, donc $\lambda = \frac{\lambda_0}{n}$

\subsection{Champ magnétique de l'onde}
On sait que $\underline{\vec{B}} = -\frac{i}{\omega}\rot \underline{\vec{E}}$ ou alors on utilise
$i\omega \underline{\vec{B}} = -\vec{k}\wedge \underline{\vec{E}}$

On obtient $\underline{\vec{B}} = \frac{n}{c}\underline{\vec{E_x}}\vec{e_y}$
\subsection{Vecteur de Poyting}
On calcule $\vec{\Pi} = \vec{E}\wedge \frac{\vec{B}}{\mu_0}$ en utilisant les vecteur réels et non complexes.

On en déduit que $\vec{\Pi} = \frac{1}{\mu_0}E_xB_y\vec{e_z}$

En outre, $I = \frac{<\dd P_r>}{\dd S} = \frac{<\vec{Pi}\cdot \vec{\dd S}>}{\dd S} = <R_z>$

On en déduit que $I = \frac12 \frac{1}{\mu_0} Re(\underline{E_x}\cdt \underline{B_y}*) = \frac12 \frac{1}{\mu_0}E_0B_0 = \frac12 \frac{1}{\mu_0}E_0^2\frac{n}{c}$

Donc 
$$I = \frac{1}{2}nc\epsilon_0E_0^2$$

Faire les app numériques

\section{L'arséniure}
On met volontairement ces impuretés.

On combine ici la polarisation naturelle et celle apportée par le dopage.

\subsection{Modèle de Drude}
On fait un bilan des forces.

On utilise le modèle de Drude : 
$$m* \frac{\dd \vec{v}}{\dd t} = \vec{F_{lorenz}}-\frac{m*}{\tau}\vec{v}$$
On néglige les collisions. Ici, $\tau$ représente le temps moyen entre 2 coliisons. Si il n'y a pas de collisison, $\tau\to \infty$, donc on néglige le terme.

Si il y avait des électrons liées, il faudrait rajouter en plus $-m\omega_0\vec{r}$.

On écrit donc 
$$m*\frac{\dd \underline{\vec{v}}}{\dd t } = -e(\underline{E}+\underline{\vec{v}}\wedge \underline{\vec{B}})$$
\subsection{Amplutude des forces}
En faisant le rapport des deux, on a 
$$\frac{\norm{F_m}}{\norm{F_e}} = \frac{\norm{v}}{c} n <<1$$
\subsection{Équation linéaire}
\begin{flalign*}
    m*\frac{\dd \underline{\vec{v}}}{\dd t}& = -\frac{-e}{m}\underline{\vec{E}}\\
   -i\omega m* \underline{\vec{v}} &= -e \underline{\vec{E}}

\end{flalign*}
Pour le courant : $$\underline{\vec{J}} = -eN\underline{\vec{v}}$$
on en déduit $\underline{\vec{J}} = i\frac{Ne^2}{m*\omega}\underline{\vec{E}}$
Comme on a une relation linéaire entre J et E, on peut introduire la conductivité $\underline{\sigma}$
\subsection{Courant de Polarisation}
Quand une polarisation varie ou cours du temps, on trouve un courant. Ici, on fait le raisonnement inverse.

On écrit : 
\begin{flalign*}
    \underline{\vec{J}} = \underline{\vec{J_{libre}}} = \frac{\partial \underline{\vec{P_{libre}}}}{\partial t}\\
    &= -i\omega \underline{\vec{P_{libre}}}\\
   \underline{\vec{P_{libre}}} &= \frac{i}{\omega} \underline{\vec{J}}\\
   &= -\frac{N_e^2}{m*\omega^2}\underline{\vec{E}}\\
   &= -\epsilon_0\epsilon_s \frac{\omega_p^2}{\omega^2}\underline{\vec{E}}
\end{flalign*}
On a introduit $\omega_p$ la pulsation plasma $\omega_p^2 = \frac{N_e^2}{m*\epsilon_0\epsilon_s}$
\subsection{Polarisation totale}
On a $\underline{\vec{P}} = \underline{\vec{P_{pur}}} + \underline{\vec{P_{libres}}}$

Donc 

\begin{flalign*}
    \underline{\vec{P}}&= \epsilon_0(\epsilon_s(1-\frac{\omega_p^2}{\omega^2})-1)\underline{\vec{E}}\\
    & = \epsilon_0\chi_e^{dope}\underline{\vec{E}}\\
    & = \epsilon_0(\epsilon_r-1)\underline{\vec{E}}
\end{flalign*}

Par identification, on trouve bien l'expression de l'énoncé. On remarque que la valeur est réelle, donc il n'y a pas d'absorpotion.
\subsection{Représentation}
En pratique, on fait varier le dopage N et donc c'est $\omega_P$ qui évolue.

On peut en déduire facilement $N_0$
\subsection{Caractère de l'onde EM}
On introduit l'indice complexe : $\underline{n} = n+i\kappa$. De plus, $\underline{k} = \underline{n}k_0$

De plus, $\underline{\epsilon_r} = \underline{n^2}$ Mais ici, c'est réel, donc on a $\epsilon_r = n^2-\kappa^2+2in\kappa$

donc $\epsilon_r = n^2-\kappa^2$ et $n\kappa=0$

Si $N<N_0$, alors $\omega_p<\omega$ et $\epsilon_r>0$ donc $\kappa=0$

On écrit donc $\underline{\vec{E}} = \underline{\vec{E_0}}e^{i(\frac{\omega}{c}nz-\omega t)}$

L'onde se propage dans le milieu sans modification de l'amplitude. On peut qualifier le milieu de transparent.

Dans le deuxième cas de figure, $\epsilon_r < 0$, on botient l'opposé : $n=0$

On écrit donc $\underline{\vec{E}} = \underline{\vec{E_0}}e^{i(\frac{\omega}{c}i\kappa z-\omega t)}$ avec une composante réelle qui décroit quand z augmente. Il n'y a pas de propagation car z et t ne sont pas couplés. 
Il apparaissent dans des fonctions différentes. C'est donc une onde évanessante.

Cet effet n'est pas en lien avec l'absorbtion. L'atténuation provient du fait que $\epsilon_r <0$ et réel ou parce qu'il est complexe.

Il faut trouver $\epsilon_s$. 

Faire un commentaire sur le m*.


\end{document}
